% !TEX root = ../main.tex

\section{Conclusion}
\label{sec:conclusion}

We asked whether differences at a small feature limit reflect search quality among small subsets or only whether the searches reach such subsets within a finite budget, and compared an opposition-based start, a sparse start, and our implementation of NSGAII-MIIP on up to 46 datasets (NSGAII-MIIP on 39), with a pilot on four and a confirmatory test on seven new datasets. At a limit of ten features, the split of the AUF difference between the two starts was dominated by reach, which made up 98 to 100 percent of the combined absolute reach and quality parts. At sizes both starts reached, the mean test F1 difference was small under two classifiers and inconclusive under the third. Reach depended on the number of features and the budget: the fits the OBL start needed to reach ten features grew as about $d^{1.7}$, a relation predicted from the pilot and confirmed on the new datasets. Under the tested variation and survival rules no start kept both ends of the front. A large AUF difference at a small limit is therefore not by itself evidence of better search, and it should be reported together with reach and the budget. Future work should test the matched-size reading on new splits and unseen datasets with the margin and the eligibility rule fixed in advance, look for a selection or archive rule that keeps both ends of the front, separate the number of features from the number of training rows in the budget relation and test it under other classifiers, and extend the comparison to wider datasets, larger budgets, and other operator families.
